% GENERATED FILE. Edit canonical lessons, then run npm run generate:books.
% canonical-source-hash: fnv1a64:0c23d3833aef3f29
% canonical-lessons: AR-C435-masar, AR-C435-ujra, AR-C435-ishtirak, AR-C435-mamarr-mushah, AR-C435-istadara

\chapter{Routes and fares}
\label{ch:ar-routes-and-fares}
\hlchaptermodality{435}

\begin{chapteropening}
\textbf{By the end of this chapter:} \emph{I can say a route, a fare, a subscription, a pedestrian crossing and to turn around.}

Complete the last lesson of chapter 435: I can say a route, a fare, a subscription, a pedestrian crossing and to turn around.
\end{chapteropening}

\section[\texorpdfstring{\ar{مسار} (masār)}{مسار (masār)}]{\ar{مسار} (masār) — a route; a lane}

\label{lesson:AR-C435-masar}



\begin{quote}
Before the new one: say the Arabic for to issue, then the Arabic for to overtake.
\end{quote}

\subsection*{You'll want to know: \ar{مسار}}
\textbf{\ar{مسار}} — \emph{masār} — \textquotedblleft{}a route; a lane\textquotedblright{}.

\begin{quote}
\ar{غيرت} \ar{الحافلة} \ar{مسارها} \ar{بسبب} \ar{الثلج}.
\end{quote}

\begin{quote}
\emph{ghayyarat al-ḥāfila masārahā bi-sabab ath-thalj.}
\end{quote}

\begin{quote}
The bus changed its route because of the snow.
\end{quote}

\begin{cousinweb}
\textbf{\ar{مسار}} (\emph{masār}) is built on the native Arabic root \emph{s-y-r}, \textquotedblleft{}to go, to travel\textquotedblright{} (Lane). The pattern \emph{mafʿal} makes it the place of going.
\end{cousinweb}

\begin{grammarlens}[title={What you've built}]
One more word for getting around.
\end{grammarlens}

\subsection*{Your turn}
\begin{itemize}
\raggedright
  \item \emph{Say it:} \emph{masār}
  \item \emph{Say it:} \emph{masār}, once more
  \item \emph{Say it:} \emph{ghayyarat al-ḥāfila masārahā bi-sabab ath-thalj.} — The bus changed its route because of the snow.
  \item \emph{Recall:} say the Arabic for to issue, then the Arabic for to overtake, then say \emph{masār} again
\end{itemize}

\begin{tcolorbox}[breakable,colback=teal!4,colframe=teal!35!black,title={Before you move on}]
What does \ar{مسار} mean? (\textquotedblleft{}A route; a lane\textquotedblright{}.) Say it once more.
\end{tcolorbox}
\section[\texorpdfstring{\ar{أجرة} (ujra)}{أجرة (ujra)}]{\ar{أجرة} (ujra) — a fare; a wage}

\label{lesson:AR-C435-ujra}



\begin{quote}
Before the new one: say the Arabic for to overtake, then the Arabic for a route.
\end{quote}

\subsection*{You'll want to know: \ar{أجرة}}
\textbf{\ar{أجرة}} — \emph{ujra} — \textquotedblleft{}a fare; a wage\textquotedblright{}.

\begin{quote}
\ar{كانت} \ar{أجرة} \ar{التاكسي} \ar{رخيصة}.
\end{quote}

\begin{quote}
\emph{kānat ujrat at-tāksī rakhīṣa.}
\end{quote}

\begin{quote}
The taxi fare was cheap.
\end{quote}

\begin{cousinweb}
\textbf{\ar{أجرة}} (\emph{ujra}) is built on the native Arabic root \emph{ʾ-j-r}, \textquotedblleft{}a wage, a reward\textquotedblright{} (Lane). A fare is the wage paid for the ride.
\end{cousinweb}

\begin{grammarlens}[title={What you've built}]
One more word for getting around.
\end{grammarlens}

\subsection*{Your turn}
\begin{itemize}
\raggedright
  \item \emph{Say it:} \emph{ujra}
  \item \emph{Say it:} \emph{ujra}, once more
  \item \emph{Say it:} \emph{kānat ujrat at-tāksī rakhīṣa.} — The taxi fare was cheap.
  \item \emph{Recall:} say the Arabic for to overtake, then the Arabic for a route, then say \emph{ujra} again
\end{itemize}

\begin{tcolorbox}[breakable,colback=teal!4,colframe=teal!35!black,title={Before you move on}]
What does \ar{أجرة} mean? (\textquotedblleft{}A fare; a wage\textquotedblright{}.) Say it once more.
\end{tcolorbox}
\section[\texorpdfstring{\ar{اشتراك} (ishtirāk)}{اشتراك (ishtirāk)}]{\ar{اشتراك} (ishtirāk) — a subscription; a travel pass}

\label{lesson:AR-C435-ishtirak}



\begin{quote}
Before the new one: say the Arabic for a route, then the Arabic for a fare.
\end{quote}

\subsection*{You'll want to know: \ar{اشتراك}}
\textbf{\ar{اشتراك}} — \emph{ishtirāk} — \textquotedblleft{}a subscription; a travel pass\textquotedblright{}.

\begin{quote}
\ar{اشترى} \ar{أخي} \ar{اشتراكا} \ar{شهريا} \ar{للمترو}.
\end{quote}

\begin{quote}
\emph{ishtarā akhī ishtirākan shahriyyan li-l-mitrū.}
\end{quote}

\begin{quote}
My brother bought a monthly pass for the metro.
\end{quote}

\begin{cousinweb}
\textbf{\ar{اشتراك}} (\emph{ishtirāk}) is built on the native Arabic root \emph{sh-r-k}, \textquotedblleft{}to share\textquotedblright{} (Lane). It is the verbal noun of Form VIII, \textquotedblleft{}taking a share in something\textquotedblright{}.
\end{cousinweb}

\begin{grammarlens}[title={What you've built}]
One more word for getting around.
\end{grammarlens}

\subsection*{Your turn}
\begin{itemize}
\raggedright
  \item \emph{Say it:} \emph{ishtirāk}
  \item \emph{Say it:} \emph{ishtirāk}, once more
  \item \emph{Say it:} \emph{ishtarā akhī ishtirākan shahriyyan li-l-mitrū.} — My brother bought a monthly pass for the metro.
  \item \emph{Recall:} say the Arabic for a route, then the Arabic for a fare, then say \emph{ishtirāk} again
\end{itemize}

\begin{tcolorbox}[breakable,colback=teal!4,colframe=teal!35!black,title={Before you move on}]
What does \ar{اشتراك} mean? (\textquotedblleft{}A subscription; a travel pass\textquotedblright{}.) Say it once more.
\end{tcolorbox}
\section[\texorpdfstring{\ar{ممر} \ar{مشاة} (mamarr mushāh)}{ممر مشاة (mamarr mushāh)}]{\ar{ممر} \ar{مشاة} (mamarr mushāh) — a pedestrian crossing}

\label{lesson:AR-C435-mamarr-mushah}



\begin{quote}
Before the new one: say the Arabic for a fare, then the Arabic for a subscription.
\end{quote}

\subsection*{You'll want to know: \ar{ممر} \ar{مشاة}}
\textbf{\ar{ممر} \ar{مشاة}} — \emph{mamarr mushāh} — \textquotedblleft{}a pedestrian crossing\textquotedblright{}.

\begin{quote}
\ar{عبر} \ar{الأطفال} \ar{الشارع} \ar{من} \ar{ممر} \ar{المشاة}.
\end{quote}

\begin{quote}
\emph{ʿabara al-aṭfāl ash-shāriʿ min mamarr al-mushāh.}
\end{quote}

\begin{quote}
The children crossed the street at the pedestrian crossing.
\end{quote}

\begin{cousinweb}
\textbf{\ar{ممر} \ar{مشاة}} (\emph{mamarr mushāh}) is built from two native Arabic words: \emph{mamarr}, \textquotedblleft{}a passage\textquotedblright{}, on the root \emph{m-r-r}, \textquotedblleft{}to pass\textquotedblright{} (Lane), and \emph{mushāh}, \textquotedblleft{}walkers\textquotedblright{}, on the root \emph{m-sh-y}, \textquotedblleft{}to walk\textquotedblright{} (Lane).
\end{cousinweb}

\begin{grammarlens}[title={What you've built}]
One more word for getting around.
\end{grammarlens}

\subsection*{Your turn}
\begin{itemize}
\raggedright
  \item \emph{Say it:} \emph{mamarr mushāh}
  \item \emph{Say it:} \emph{mamarr mushāh}, once more
  \item \emph{Say it:} \emph{ʿabara al-aṭfāl ash-shāriʿ min mamarr al-mushāh.} — The children crossed the street at the pedestrian crossing.
  \item \emph{Recall:} say the Arabic for a fare, then the Arabic for a subscription, then say \emph{mamarr mushāh} again
\end{itemize}

\begin{tcolorbox}[breakable,colback=teal!4,colframe=teal!35!black,title={Before you move on}]
What does \ar{ممر} \ar{مشاة} mean? (\textquotedblleft{}A pedestrian crossing\textquotedblright{}.) Say it once more.
\end{tcolorbox}
\section[\texorpdfstring{\ar{استدار} (istadāra)}{استدار (istadāra)}]{\ar{استدار} (istadāra) — to turn around}

\label{lesson:AR-C435-istadara}



\begin{quote}
Before the new one: say the Arabic for a subscription, then the Arabic for a pedestrian crossing.
\end{quote}

\subsection*{You'll want to know: \ar{استدار}}
\textbf{\ar{استدار}} — \emph{istadāra} — \textquotedblleft{}to turn around\textquotedblright{}.

\begin{quote}
\ar{استدار} \ar{الولد} \ar{عندما} \ar{سمع} \ar{اسمه}.
\end{quote}

\begin{quote}
\emph{istadāra al-walad ʿindamā samiʿa ismahu.}
\end{quote}

\begin{quote}
The boy turned around when he heard his name.
\end{quote}

Present: \ar{يستدير}, \emph{yastadīr}.

\begin{cousinweb}
\textbf{\ar{استدار}} (\emph{istadāra}) is built on the native Arabic root \emph{d-w-r}, \textquotedblleft{}to turn, to go round\textquotedblright{} (Lane). It is Form X.
\end{cousinweb}

\begin{grammarlens}[title={What you've built}]
One more word for getting around.
\end{grammarlens}

\subsection*{Your turn}
\begin{itemize}
\raggedright
  \item \emph{Say it:} \emph{istadāra}
  \item \emph{Say it:} \emph{istadāra}, once more
  \item \emph{Say it:} \emph{istadāra al-walad ʿindamā samiʿa ismahu.} — The boy turned around when he heard his name.
  \item \emph{Recall:} say the Arabic for a subscription, then the Arabic for a pedestrian crossing, then say \emph{istadāra} again
\end{itemize}

\begin{tcolorbox}[breakable,colback=teal!4,colframe=teal!35!black,title={Before you move on}]
What does \ar{استدار} mean? (\textquotedblleft{}To turn around\textquotedblright{}.) Say it once more.
\end{tcolorbox}
